% ============================================================
% Resumen y English Abstract
% ============================================================

\chapter*{Resumen}
\addcontentsline{toc}{chapter}{Resumen}
\markright{Resumen}
\enlargethispage{3\baselineskip}

Este trabajo estudia la topología computacional y su uso en la
segmentación de imágenes. Se exponen los fundamentos de la homología persistente:
complejos simpliciales y cúbicos, grupos de homología y números de
Betti, filtraciones y diagramas de persistencia, su cálculo por
reducción matricial o, en dimensión $0$, por \emph{union-find}, y su
estabilidad frente a perturbaciones. Después se
presentan los métodos que llevan la topología a la segmentación, en
particular Topograph, una pérdida que penaliza solo las regiones que
cambian la topología de la predicción.

Estos conocimientos se aplican a la segmentación de filamentos solares
en imágenes H-alfa, un problema en el que los errores que más pesan son
topológicos: filamentos partidos, fusionados o manchas espurias. La red asigna a
cada píxel una probabilidad de ser filamento, y el post-proceso
habitual conserva los píxeles que superan un umbral y descarta las
componentes pequeñas. Al bajar el umbral, cada componente nace en el
valor de su máximo y muere al unirse con otra más alta, y ese par de
valores es un punto del diagrama de persistencia $H_0$. Se demuestra
que las instancias que deja un umbral son los puntos nacidos por encima
de él y aún vivos, así que fijar el umbral es elegir qué puntos del
diagrama pasan a ser instancias. Cada error se reconoce
además en el diagrama. Un filamento partido da varios puntos, dos
fusionados dejan uno solo y una mancha espuria es un punto sin
filamento debajo. Esta lectura
guía dos formas de corregir esos errores: durante el entrenamiento, con
pérdidas topológicas como clDice y Topograph, y después, con un
post-proceso que decide qué componentes conservar según su persistencia.

Topograph hace predicciones más fiables, con mayor precisión y muchas
menos predicciones espurias, aunque en la métrica se obtienen
resultados similares simplemente filtrando los segmentos por área. El post-proceso por
persistencia conserva la estructura y reduce la fragmentación, pero no
separa las predicciones espurias muy seguras, porque ahí la
persistencia apenas se distingue de la confianza del modelo. Combinado con el área y con
Topograph, da los menores errores de estructura.
En conjunto, la topología aporta
una forma precisa de entender y clasificar los errores de la
segmentación, y criterios para corregirlos que el tamaño no
ofrece.

\chapter*{English Abstract}
\addcontentsline{toc}{chapter}{English Abstract}
\markright{English Abstract}
\enlargethispage{3\baselineskip}

\begin{otherlanguage}{english}
This work studies computational topology and its use in image
segmentation. It presents the foundations of persistent homology:
simplicial and cubical complexes, homology groups and Betti numbers,
filtrations and persistence diagrams, their computation by matrix
reduction or, in dimension $0$, by union-find, and their stability
under perturbations. It then presents
the methods
that bring topology into segmentation, in particular Topograph, a loss
that penalises only the regions that change the topology of the
prediction.

This knowledge is applied to the segmentation of solar filaments in
H-alpha images, a problem in which the most costly errors are
topological: split or merged filaments and spurious blobs. The network assigns
each pixel a probability of being a filament, and the usual
post-processing keeps the pixels above a threshold and discards small
components. As the threshold is lowered, each component is born at the
value of its maximum and dies when it merges with a higher one, and
that pair of values is a point of the $H_0$ persistence diagram. We
prove that the instances left by a threshold are the points born above
it and still alive, so fixing the threshold amounts to choosing which
points of the diagram become instances. Each error can also be
recognised in the diagram. A split filament gives several points, two
merged filaments leave only one, and a spurious blob is a point with no
filament underneath. This reading guides two ways of
correcting those errors: during training, with topological losses such
as clDice and Topograph, and afterwards, with a post-processing that
decides which components to keep according to their persistence.

Topograph produces more reliable predictions, with higher precision and
far fewer spurious predictions, although similar results in the metric
are obtained simply by filtering the segments by area. Persistence-based
post-processing preserves structure and reduces fragmentation, but does
not separate highly confident spurious predictions, because there
persistence barely differs from the model's confidence. Combined with the
area filter and with Topograph, it yields the fewest structural errors. Overall, topology provides a precise
way to understand and classify segmentation errors, and criteria to
correct them that component size alone does not offer.
\end{otherlanguage}
